This thesis uses observational data primarily from the ZFOURGE Survey in addition to other sources that will be discussed below. It is therefore important that the benefits and constraints that are imposed by using these datasets.

% Refer to the Straatman paper and Tomzcak paper to further detail the information to go into this section
\section{ZFOURGE}
The data that is used in this thesis is from the the FourStar Galaxy Evolution Survey (ZFOURGE) \citep{straatman_fourstar_2016}, a large photometric imaging survey covering approximately $400$ arcmin$^2$ over three of the Hubble Space Telescope (HST) legacy fields: CDFS \citep{giacconi_chandra_2002}, COSMOS \citep{scoville_cosmic_2007}, and UDS \citep{lawrence_ukirt_2007}. To collect observational data, the ZFOURGE Survey employed the use of the FourStar Infrared Camera, which was mounted on the Magellan Baade Telescope at the Las Campanas Observatory in Chile \citep{persson_fourstar_2013}. Being a photometric survey, ZFOURGE measures the light of distant galaxies using different band-pass filters. While many surveys use broadband filters, the strength of ZFOURGE lies in the fact that it uses deep near-IR imaging via the $J_1$, $J_2$, $J_3$ $H_1$, \& $H_2$ medium band filters and the $K_s$ broadband filter to provide highly reliable photometric redshift measurements within $\Delta z/(1+z) \simeq 1 -2\%$ at $1 < z < 4$ \citep{straatman_fourstar_2016}. In addition to this, ZFOURGE incorporates multi-wavelength data from Cosmic Assembly Near-infrared Deep Extragalactic Legacy Survey (CANDELS) \citep{koekemoer_candels_2011, grogin_candels_2011, skelton_3d-hst_2014}, alongside imaging from Spitzer\cite{fazio_infrared_2004, rieke_multiband_2004} and other ground-based telescopes. The Rest-frame colours and photometric redshifts in ZFOURGE were calculated using EAZY \citep{brammer_eazy_2008}. The data, reduction methods and photometric bands can be fully described in \cite{straatman_fourstar_2016}.


% Present a larger discussion of the supplementary sources
\section{Supplementary Sources}
While the ZFOURGE survey contains a wealth of data in UV, optical and IR wavelengths, ZFOURGE is also supplemented with multi-wavelength data in the radio, x-ray and far-IR to ensure that AGN can be investigated at all wavelengths and from all angles \citep{padovani_active_2017}. To enhance its depth and breadth, the ZFOURGE dataset is further augmented with complementary multiwavelength observations from both the Herschel Space Observatory, the Chandra X-ray Observatory and the XMM-Newton space telescope \citep{cowley_zfourge_2016}.

